\documentclass[11pt,a4paper]{ctexart}
\usepackage[margin=2.2cm]{geometry}
\usepackage{amsmath,amssymb}
\usepackage{fancyhdr}\usepackage{enumitem}\usepackage{xcolor}
\usepackage{hyperref}
\pagestyle{fancy}\fancyhf{}
\rhead{Calculus \& Linear Algebra}
\lhead{Songpengfei}
\rfoot{Page \thepage}
\newcommand{\solvertag}[1]{{\footnotesize\color{gray}\textsf{[#1]}}}

\title{Homework 3 Solutions}
\author{Songpengfei}
\date{2026-10-09}
\begin{document}
\maketitle

\noindent\textbf{Solved: 10/10 (100.0\%)}\quad\textit{Auto-generated by C4C homework-solver.}

\section*{Problem PartA1}
Find the derivative of $x^3 + 2x^2 - 5x + 1$.

\paragraph{Solution.} \solvertag{sympy.calculus} 

\begin{enumerate}[leftmargin=*]
  \item 已知 $f(x) = x^{3} + 2 x^{2} - 5 x + 1$。
  \item 求一阶导：$f'(x) = \dfrac{d}{dx}\left(x^{3} + 2 x^{2} - 5 x + 1\right)$。
  \item 逐项使用 $(x^n)' = nx^{n-1}$、$(\sin x)' = \cos x$、$(e^x)' = e^x$、$(\ln x)' = 1/x$ 与乘积/链式法则。
  \item 得 $f'(x) = 3 x^{2} + 4 x - 5$。
\end{enumerate}

\noindent\textbf{Answer:}\quad $3 x^{2} + 4 x - 5$

\medskip

\section*{Problem PartA2}
Evaluate the integral $\int x^2 \sin(x) dx$.

\paragraph{Solution.} \solvertag{sympy.calculus} 

\begin{enumerate}[leftmargin=*]
  \item 求不定积分 $\displaystyle\int x^{2} \sin{\left(x \right)}\, dx$。
  \item 选择合适的技巧（换元 / 分部积分 / 部分分式 / 基本积分表）。
  \item 得 $F(x) = - x^{2} \cos{\left(x \right)} + 2 x \sin{\left(x \right)} + 2 \cos{\left(x \right)} + C$。
  \item 核验：$\dfrac{d}{dx}\left(- x^{2} \cos{\left(x \right)} + 2 x \sin{\left(x \right)} + 2 \cos{\left(x \right)}\right) = x^{2} \sin{\left(x \right)} = $ 被积函数 ✓。
\end{enumerate}

\noindent\textbf{Answer:}\quad $C - x^{2} \cos{\left(x \right)} + 2 x \sin{\left(x \right)} + 2 \cos{\left(x \right)}$

\medskip

\section*{Problem PartA3}
Compute the limit $\lim_{x \to 0} \frac{\sin(x)}{x}$.

\paragraph{Solution.} \solvertag{sympy} 

\begin{enumerate}[leftmargin=*]
  \item $\lim_{x \to 0} \frac{\sin{\left(x \right)}}{x} = 1$
\end{enumerate}

\noindent\textbf{Answer:}\quad $1$

\medskip

\section*{Problem PartA4}
Find the derivative of $e^{x^2} \cdot \ln(x)$.

\paragraph{Solution.} \solvertag{sympy.calculus} 

\begin{enumerate}[leftmargin=*]
  \item 已知 $f(x) = e^{x^{2}} \log{\left(x \right)}$。
  \item 求一阶导：$f'(x) = \dfrac{d}{dx}\left(e^{x^{2}} \log{\left(x \right)}\right)$。
  \item 逐项使用 $(x^n)' = nx^{n-1}$、$(\sin x)' = \cos x$、$(e^x)' = e^x$、$(\ln x)' = 1/x$ 与乘积/链式法则。
  \item 得 $f'(x) = \frac{\left(2 x^{2} \log{\left(x \right)} + 1\right) e^{x^{2}}}{x}$。
\end{enumerate}

\noindent\textbf{Answer:}\quad $\frac{\left(2 x^{2} \log{\left(x \right)} + 1\right) e^{x^{2}}}{x}$

\medskip

\section*{Problem PartA5}
Evaluate $\int_0^1 \frac{1}{1+x^2} dx$.

\paragraph{Solution.} \solvertag{sympy.calculus} 

\begin{enumerate}[leftmargin=*]
  \item 求定积分 $\displaystyle\int_{0}^{1} \frac{1}{x^{2} + 1}\, dx$。
  \item 先求原函数 $F(x) = \operatorname{atan}{\left(x \right)}$。
  \item 牛顿-莱布尼茨公式：$F(1) - F(0) = \frac{\pi}{4}$。
\end{enumerate}

\noindent\textbf{Answer:}\quad $\frac{\pi}{4}$

\medskip

\section*{Problem PartB6}
Solve the equation $x^2 - 5x + 6 = 0$.

\paragraph{Solution.} \solvertag{sympy} 

\begin{enumerate}[leftmargin=*]
  \item $x^{2} - 5 x + 6 = 0$
  \item $x = \left[ 2, \  3\right]$
\end{enumerate}

\noindent\textbf{Answer:}\quad $\left[ 2, \  3\right]$

\medskip

\section*{Problem PartB7}
Solve $x^3 - 6x^2 + 11x - 6 = 0$.

\paragraph{Solution.} \solvertag{sympy} 

\begin{enumerate}[leftmargin=*]
  \item $x^{3} - 6 x^{2} + 11 x - 6 = 0$
  \item $x = \left[ 1, \  2, \  3\right]$
\end{enumerate}

\noindent\textbf{Answer:}\quad $\left[ 1, \  2, \  3\right]$

\medskip

\section*{Problem PartC8}
Compute the determinant of the matrix $A = \begin{pmatrix} 1 & 2 \\ 3 & 4 \end{pmatrix}$.

\paragraph{Solution.} \solvertag{sympy.linalg} 

\begin{enumerate}[leftmargin=*]
  \item 取矩阵 $A = \left[\begin{matrix}1 & 2\\3 & 4\end{matrix}\right]$，阶数 $n = 2$。
  \item 按定义展开：$\det A = \sum_{\sigma \in S_n} \operatorname{sgn}(\sigma)\prod_{i=1}^{n} a_{i,\sigma(i)}$（$n=2$ 时共 $2$ 项）。
  \item 用余子式展开/行变换化简，得 $\det A = -2$。
  \item 核对：$\det A = -2= -2$。
\end{enumerate}

\noindent\textbf{Answer:}\quad $-2$

\medskip

\section*{Problem PartC9}
Find the eigenvalues of $A = \begin{pmatrix} 2 & 1 \\ 1 & 2 \end{pmatrix}$.

\paragraph{Solution.} \solvertag{sympy.linalg} 

\begin{enumerate}[leftmargin=*]
  \item $A = \left[\begin{matrix}2 & 1\\1 & 2\end{matrix}\right]$。
  \item 特征方程 $\det(A - \lambda I) = 0$，即 $\lambda^{2} - 4 \lambda + 3 = 0$。
  \item 解得特征值 $\lambda = 1$。
  \item 解得特征值 $\lambda = 3$。
\end{enumerate}

\noindent\textbf{Answer:}\quad $\lambda = 1,\ 3$

\medskip

\section*{Problem PartD10}
Solve the ODE $y' + 2y = e^{-x}$.

\paragraph{Solution.} \solvertag{sympy.ode} 

\begin{enumerate}[leftmargin=*]
  \item 原方程：$2 y{\left(x \right)} + \frac{d}{d x} y{\left(x \right)} = e^{- x}$。
  \item 这是一个 1 阶线性常微分方程。
  \item 解得：$$y{\left(x \right)} = \left(C_{1} e^{- x} + 1\right) e^{- x}$$
  \item （通解中的 $C_1, C_2, \dots$ 为任意常数；给出初值即可确定。）
  \item 核验：把解代回原方程，两边恒等 ✓。
\end{enumerate}

\noindent\textbf{Answer:}\quad $\left(C_{1} e^{- x} + 1\right) e^{- x}$

\medskip

\end{document}
